% ==============================================================================
% Section 2: Channel Mixer (Feed-Forward Networks)
% Full LaTeX body — 8 subsections, no \section command
% ==============================================================================

\subsection{FFN Structure and Activation Functions}
\label{subsec:ffn_activation}

The feed-forward network (FFN) within each Transformer block serves as the channel mixer, applying an identical pointwise nonlinear transformation to every token representation independently.
In the original Transformer architecture~\cite{vaswani2017attention}, the FFN is a two-layer perceptron with a single hidden expansion:
\begin{equation}
\label{eq:vanilla_ffn}
\mathrm{FFN}(x) = W_2 \,\sigma\!\bigl(W_1 x + b_1\bigr) + b_2,
\end{equation}
where $W_1 \in \mathbb{R}^{d \times d_{\mathrm{ff}}}$, $W_2 \in \mathbb{R}^{d_{\mathrm{ff}} \times d}$, and $\sigma$ denotes a pointwise activation function.
The hidden dimension $d_{\mathrm{ff}}$ is typically set to $4d$, so the FFN accounts for roughly two-thirds of each block's parameters.
Because the activation function is the sole source of nonlinearity in this module, its choice fundamentally determines the expressiveness of the per-token feature transformation.

The earliest Transformer models adopted \textbf{ReLU}~\cite{agarap2018deep}, $\mathrm{ReLU}(x)=\max(0,x)$, which is computationally efficient but suffers from the dying neuron problem: units whose pre-activations are consistently negative receive zero gradient and cease to learn.
\textbf{GELU}~\cite{hendrycks2016gaussian} addresses this shortcoming by interpreting activation as a stochastic regularisation expectation, $\mathrm{GELU}(x)=x\,\Phi(x)$, where $\Phi$ is the standard Gaussian CDF.
Its smoothness yields a more favourable optimisation landscape, and GELU became the default in GPT-2 and BERT.
\textbf{Swish} (also called SiLU)~\cite{ramachandran2017searching}, discovered via automated activation search, generalises the sigmoid gate to $\mathrm{Swish}(x)=x\,\sigma(\beta x)$; with $\beta=1$ it coincides with SiLU.
Unlike ReLU, Swish is non-monotonic and everywhere differentiable, properties that empirically benefit deep networks.

A more recent line of work re-examines ReLU-family activations from a sparsity perspective.
\textbf{ReLU$^2$}~\cite{zhang2024relu2} squares the ReLU output, $\mathrm{ReLU}^2(x)=\bigl(\max(0,x)\bigr)^2$, which amplifies the zero region and produces activation sparsity exceeding 90\%.
This hard sparsity---where inactive neurons are exactly zero rather than merely small---enables hardware-level shortcuts that skip zero-valued multiplications, achieving inference speedups comparable to or better than dense SwiGLU networks with negligible quality loss.
The tension between smooth activations that ease training and sparse activations that accelerate inference remains an active area of investigation; emerging approaches such as \textbf{PolyGLU} attempt to reconcile the two by conditioning the activation mode on the input state, breaking the static one-activation-fits-all paradigm.


% ------------------------------------------------------------------------------
\subsection{Gated Linear Units and the Three-Matrix FFN}
\label{subsec:glu}

The introduction of \textbf{Gated Linear Units} (GLU)~\cite{dauphin2017language} marked a structural turning point for the Transformer FFN, replacing the standard two-matrix architecture with a three-matrix design that incorporates multiplicative gating.
In the original formulation, GLU computes
\begin{equation}
\label{eq:glu}
\mathrm{GLU}(x) = \bigl(x W + b\bigr) \otimes \sigma\!\bigl(x V + c\bigr),
\end{equation}
where $\sigma$ is the sigmoid function, $\otimes$ denotes element-wise multiplication, and the right branch acts as a gate that controls the information flow of the left branch.
This mechanism is conceptually analogous to the forget and input gates of LSTMs: one projection generates candidate features while a parallel projection determines which features are retained.

Shazeer~\cite{shazeer2020glu} systematically evaluated GLU variants by substituting the sigmoid gate with different activation functions.
\textbf{ReGLU} uses ReLU, \textbf{GeGLU} uses GELU, and \textbf{SwiGLU} uses Swish:
\begin{equation}
\label{eq:swiglu}
\mathrm{SwiGLU}(x) = \mathrm{Swish}\!\bigl(x W_{\mathrm{gate}}\bigr) \otimes \bigl(x W_{\mathrm{up}}\bigr), \qquad
\mathrm{FFN}_{\mathrm{SwiGLU}}(x) = \bigl(\mathrm{SwiGLU}(x)\bigr) W_{\mathrm{down}}.
\end{equation}
Across matched-parameter comparisons, SwiGLU consistently outperformed both the vanilla FFN and other GLU variants, establishing it as the de facto standard in modern LLMs including LLaMA~\cite{touvron2023llama}, PaLM, and Qwen.
GeGLU, which replaces Swish with GELU in the gate branch, is adopted by models such as Gemma~\cite{gemma2024improving}:
\begin{equation}
\mathrm{GeGLU}(x) = \left(\mathrm{GELU}(x W_1) \odot x V\right) W_2
\label{eq:geglu}
\end{equation}

Because the three-matrix structure introduces an additional projection ($W_{\mathrm{gate}}$, $W_{\mathrm{up}}$, $W_{\mathrm{down}}$ versus the original $W_1$, $W_2$), it increases the raw parameter count from $2 d \cdot d_{\mathrm{ff}}$ to $3 d \cdot d_{\mathrm{ff}}$.
In practice, the hidden dimension is reduced to $d_{\mathrm{ff}} = \tfrac{2}{3} \times 4d \approx \tfrac{8}{3}d$ so that the total parameter budget remains unchanged.
Even under this iso-parameter constraint, the GLU-based FFN consistently surpasses the wider two-matrix FFN, suggesting that the multiplicative interaction provides a qualitatively different---and superior---form of expressiveness.
Theoretically, the gating branch supplies a data-dependent soft mask over the feature space, enabling the network to learn higher-order feature interactions that are difficult to express with additive activations alone.
The gradient flow through the element-wise product also provides a more direct path reminiscent of residual connections, potentially easing optimisation in very deep models.

The three-matrix GLU FFN has become so dominant that the two-matrix design is essentially absent from frontier models released after 2023.
Looking ahead, the gating principle is beginning to diffuse into other modules---gated attention~\cite{qiu2025gated} applies a similar multiplicative control to attention weights---and the interplay between FFN gating and MoE routing (Section~\ref{subsec:moe_routing}) introduces a dual-gating dynamic whose theoretical implications remain largely unexplored.


% ------------------------------------------------------------------------------
\subsection{Mixture-of-Experts: Foundations and Routing Strategies}
\label{subsec:moe_routing}

The Mixture-of-Experts (MoE) paradigm replaces the single dense FFN with a collection of $N$ expert sub-networks, of which only a small subset is activated for each token, thereby decoupling model capacity from per-token computation cost~\cite{jacobs1991adaptive,cai2024survey}.
A lightweight router (or gating network) $g$ maps each token representation $x$ to a distribution over experts; the output is the weighted sum of the selected experts' outputs:
\begin{equation}
\label{eq:moe}
\mathrm{MoE}(x) = \sum_{i=1}^{N} g_i(x)\, E_i(x), \qquad
g(x) = \mathrm{TopK}\!\bigl(\mathrm{softmax}(W_r x + \epsilon),\, k\bigr),
\end{equation}
where $E_i$ denotes the $i$-th expert FFN, $W_r$ is the router's weight matrix, and $\epsilon$ is optional noise injected during training to encourage exploration.

\textbf{Top-$k$ routing}~\cite{shazeer2017outrageously} is the foundational strategy: each token selects the $k$ highest-scoring experts and discards the rest.
The original formulation uses $k{=}2$ with an auxiliary load-balancing loss that penalises uneven expert utilisation.
\textbf{Switch Transformer}~\cite{fedus2022switch} simplified this to $k{=}1$, halving communication overhead but increasing the risk of load imbalance and routing collapse---a failure mode in which a small subset of experts attracts the majority of tokens while others remain underutilised.
To stabilise training, \textbf{ST-MoE}~\cite{zoph2022stmoe} introduced the router z-loss, which penalises large logit magnitudes in the router to prevent numerical instability.

A conceptual inversion of the routing direction was proposed by \textbf{Expert Choice}~\cite{zhou2022mixture}: instead of tokens selecting experts, each expert selects a fixed number of tokens from the global pool.
This guarantees perfect load balance by construction but introduces a new asymmetry---some tokens may be processed by many experts while others are selected by none, potentially causing information loss.
\begin{equation}
S = \mathrm{softmax}(X W_g)^\top \in \mathbb{R}^{E \times N}, \quad \mathcal{T}_e = \mathrm{top\text{-}}k\!\left(S_{e,:}\right) \quad \text{(each expert selects its top-}k\text{ tokens)}
\label{eq:expert_choice}
\end{equation}

More recent work seeks to eliminate the auxiliary loss entirely.
\textbf{Auxiliary-loss-free routing}~\cite{wang2024auxiliary}, adopted by DeepSeek-V3~\cite{liu2024deepseek}, introduces a learnable bias term in the router that dynamically steers tokens from overloaded experts to underutilised ones.
By removing the auxiliary loss, this approach avoids gradient interference with the primary language-modelling objective and simplifies hyper-parameter tuning.
\begin{equation}
b_e \leftarrow b_e + \eta \cdot (\bar{f} - f_e), \quad g(x) = \mathrm{top\text{-}}k\!\left(\mathrm{softmax}(x W_g + b)\right) \quad \text{(bias-corrected routing)}
\label{eq:auxfree}
\end{equation}
\textbf{Soft MoE}~\cite{puigcerver2024soft} takes a different path by eliminating discrete routing altogether: every expert processes a soft weighted combination of all tokens, making the entire module fully differentiable at the cost of losing the computational savings of sparsity.
\textbf{Sigma-MoE}~\cite{hu2024sigma} replaces softmax with per-expert sigmoid scoring, decoupling expert weights from inter-expert competition and improving expert utilisation.
\textbf{ReMoE}~\cite{wang2024remoe} substitutes ReLU for softmax in the router, naturally producing sparse routing weights without an explicit top-$k$ selection step.

Table~\ref{tab:routing_comparison} summarises the key properties of these routing strategies.

\begin{table}[t]
\centering
\caption{Comparison of MoE routing strategies. ``Discrete'' indicates hard token-to-expert assignment; ``Load Bal.'' denotes the mechanism used to encourage balanced utilisation; ``Diff.'' indicates full differentiability.}
\label{tab:routing_comparison}
\small
\begin{tabular}{lccccl}
\hline
\textbf{Strategy} & \textbf{Direction} & \textbf{Discrete} & \textbf{Load Bal.} & \textbf{Diff.} & \textbf{Representative Model} \\
\hline
Top-$k$ + Aux Loss   & Token$\to$Expert & Yes & Aux loss       & No  & GShard~\cite{lepikhin2020gshard} \\
Top-1 (Switch)        & Token$\to$Expert & Yes & Aux loss       & No  & Switch Transformer~\cite{fedus2022switch} \\
Router z-loss         & Token$\to$Expert & Yes & z-loss         & No  & ST-MoE~\cite{zoph2022stmoe} \\
Expert Choice         & Expert$\to$Token & Yes & By design      & No  & EC-MoE~\cite{zhou2022mixture} \\
Aux-Loss-Free Bias    & Token$\to$Expert & Yes & Learned bias   & No  & DeepSeek-V3~\cite{liu2024deepseek} \\
Soft MoE              & Soft (all-to-all) & No  & N/A           & Yes & Soft MoE~\cite{puigcerver2024soft} \\
Sigma-MoE             & Token$\to$Expert & Yes & Sigmoid (ind.) & No  & Sigma-MoE~\cite{hu2024sigma} \\
ReMoE                 & Token$\to$Expert & Yes & ReLU sparsity  & No  & ReMoE~\cite{wang2024remoe} \\
\hline
\end{tabular}
\end{table}


% ------------------------------------------------------------------------------
\subsection{MoE Expert Design: Granularity, Sharing, and Heterogeneity}
\label{subsec:moe_expert_design}

Beyond routing, the design of the experts themselves---their size, number, and structural uniformity---is a critical architectural decision that governs the capacity--efficiency trade-off in MoE models.

\paragraph{Coarse- versus fine-grained experts.}
Early MoE language models employ a coarse-grained design with a small number of full-sized experts.
Mixtral~\cite{jiang2024mixtral} uses 8 experts with top-2 routing, yielding 47B total parameters of which only 13B are active per token.
This design is straightforward to implement and incurs modest communication overhead, but limits the combinatorial diversity of expert combinations.
\textbf{Fine-grained experts}, pioneered by DeepSeekMoE~\cite{dai2024deepseekmoe}, decompose each expert into a smaller FFN and dramatically increase the expert count---DeepSeek-V2~\cite{deepseekv2_2024} uses 160 routed experts with top-6 activation, while DBRX~\cite{databricks2024dbrx} employs 16 experts with top-4 selection.
The combinatorial advantage is substantial: selecting $k$ out of $N$ fine-grained experts yields $\binom{N}{k}$ possible expert combinations, providing far richer functional coverage than a coarse-grained alternative with the same active parameter count.
Pushing this idea to the extreme, He et al.~\cite{he2024million} explore million-scale experts where each expert is reduced to a minimal parameter group, paired with product-key retrieval for efficient routing.

\paragraph{Shared experts.}
DeepSeekMoE introduced the concept of \textbf{shared experts}---one or two experts that are always activated regardless of routing decisions---to capture universal, input-agnostic features common across all tokens.
The routed experts then specialise in domain- or task-specific patterns, mirroring a ``common-mode plus differential'' signal decomposition.
This design was retained in DeepSeek-V2 and V3 and has been influential in the community.
However, Qwen3~\cite{qwen2025qwen3} explicitly removes shared experts, arguing that under a global batch-level load-balancing strategy, routed experts can naturally learn general-purpose representations without a dedicated shared component.
This architectural disagreement---shared versus no-shared---reflects an unresolved question about the optimal division of labour between general and specialised computation.

\paragraph{Heterogeneous experts.}
The standard assumption that all experts share the same architecture is challenged by \textbf{heterogeneous MoE} designs such as ERNIE~4.5~\cite{ernie45_2025}, which assigns different hidden dimensions or depths to different experts.
This allows the model to allocate varying computational budgets to different input patterns, aligning with the intuition that not all tokens require the same complexity of processing.
Heterogeneous designs raise new questions for routing: the router must implicitly learn a mapping between token difficulty and expert capacity, a correspondence that is non-trivial to supervise.
Llama~4~\cite{meta2025llama4} also adopts an MoE architecture at scale, contributing to the growing diversity of expert design philosophies in production systems.


% ------------------------------------------------------------------------------
\subsection{MoE Engineering: Distributed Training and Inference}
\label{subsec:moe_engineering}

Translating the theoretical efficiency of MoE into practical speedups at scale requires addressing a constellation of systems-level challenges, from communication bottlenecks to numerical stability.

\paragraph{Expert parallelism and communication.}
\textbf{Expert Parallelism (EP)}~\cite{lepikhin2020gshard} distributes experts across GPUs so that each device holds a subset of the expert pool.
During forward propagation, an All-to-All communication primitive dispatches each token to the GPU hosting its assigned expert and returns the result.
The communication volume scales as $O(N \cdot d / P)$, where $N$ is the token count, $d$ is the hidden dimension, and $P$ is the number of devices; at large cluster sizes, this overhead can dominate computation.
\textbf{DeepEP}~\cite{deepseek2025deepep} is a purpose-built EP communication library that decomposes the All-to-All into pipelined dispatch and combine stages overlapped with expert computation, exploits NVLink topology-aware routing to minimise cross-node traffic, and supports FP8 communication to halve bandwidth requirements.

\paragraph{Load balancing in practice.}
Beyond routing-level strategies (Section~\ref{subsec:moe_routing}), systems-level load balancing involves setting a capacity factor that caps the maximum number of tokens each expert can process.
Tokens exceeding this capacity are either dropped---incurring information loss---or re-routed to alternative experts.
Switch Transformer~\cite{fedus2022switch} demonstrated that capacity factors of 1.0--1.25 provide a good balance between load uniformity and token coverage.
The auxiliary-loss-free bias approach of DeepSeek-V3 operates at the system level by dynamically adjusting per-expert biases based on running load statistics, achieving near-uniform utilisation without auxiliary gradients.

\paragraph{Mixed-precision training.}
DeepSeek-V3~\cite{liu2024deepseek} achieved a milestone by successfully training a large-scale MoE model with \textbf{FP8 mixed precision} for the expert FFN GEMMs while retaining BF16 for attention and routing computations.
The key enabler is per-group scaling---computing scale factors at a fine granularity (e.g., per 128-element tile) rather than per-tensor---to accommodate the limited dynamic range of E4M3 (maximum representable value $\approx$240).
This approach reduces communication bandwidth by approximately 40\% and accelerates expert computation, with negligible quality degradation.

\paragraph{Training stability.}
MoE training is notoriously prone to instability.
ST-MoE~\cite{zoph2022stmoe} identified that unbounded router logits can trigger gradient explosions and introduced the router z-loss $\mathcal{L}_z = \frac{1}{B}\sum_{i=1}^{B}\bigl(\log \sum_{j=1}^{N} e^{z_{ij}}\bigr)^2$ to regularise logit magnitudes.
Skywork-MoE~\cite{wei2024skywork} and OLMoE~\cite{muennighoff2024olmoe} further document best practices for learning-rate scheduling, expert initialisation, and gradient clipping that are critical for stable convergence.

\paragraph{The inference double penalty.}
While MoE reduces per-token FLOPs during training, inference presents a double penalty~\cite{moe_double_penalty_2026}: (i) all expert parameters must reside in memory even though only a fraction is used per token, resulting in extremely low parameter utilisation; and (ii) distributed inference incurs All-to-All latency proportional to expert count and cluster size.
Consequently, the training efficiency advantage of MoE does not automatically translate to inference efficiency, and bridging this gap---through expert offloading, predictive caching, or specialised hardware---remains one of the most pressing open problems.


% ------------------------------------------------------------------------------
\subsection{Conditional Computation: Beyond Expert Selection}
\label{subsec:conditional_computation}

MoE addresses the question of which experts to activate, but conditional computation asks a more general question: for a given input, how much computation is necessary at all?
This broader framework, first articulated by Bengio~\cite{bengio2013estimating}, posits that large neural networks contain far more capacity than any single input requires, and that learning to dynamically select the appropriate computational sub-graph can decouple model capacity from inference cost without sacrificing quality.

\paragraph{Theoretical foundations.}
Bengio's conditional computation framework rests on three premises: (i)~model capacity and computational cost can be decoupled through sparse activation; (ii)~the routing decision itself must be lightweight, lest the overhead of deciding what to compute negate the savings; and (iii)~the routing policy can be learned end-to-end, potentially via reinforcement learning or straight-through estimators.
The framework generalises MoE---which conditions on expert identity---to conditioning on depth (which layers to execute), width (how many channels to activate), or precision (at what numerical resolution to compute).

\paragraph{Mixture-of-Depths.}
\textbf{Mixture-of-Depths (MoD)}~\cite{raposo2024mixture} operationalises depth-wise conditional computation by training a per-layer binary router that decides, for each token, whether to execute the full Transformer sub-block (attention plus FFN) or to skip it via the residual connection.
A layer-level capacity budget (e.g., 50\% of tokens) constrains the router so that only the most ``informative'' tokens receive full computation.
Empirically, MoD achieves comparable perplexity to the dense baseline while reducing total FLOPs by more than 50\%, validating the hypothesis that natural language tokens exhibit highly non-uniform computational demand---function words such as ``the'' and ``is'' can be processed cheaply, while semantically dense tokens (negations, rare entities) require deeper processing.

\paragraph{Adaptive computation frameworks.}
\textbf{Duo-LLM}~\cite{duollm2024} maintains two parallel FFN pathways per layer---a lightweight and a full-capacity module---and routes each token to the appropriate pathway based on a learned difficulty estimate.
This design systematises the study of per-token compute allocation and can be viewed as a continuous relaxation of the binary skip/compute decision in MoD.
The concept connects to Graves' Adaptive Computation Time (ACT), which allows Universal Transformers to iteratively refine token representations with a variable number of processing steps.

\paragraph{Towards joint depth--width conditioning.}
The combination of MoE (width-level sparsity) and MoD (depth-level sparsity) yields a doubly sparse architecture in which both the set of active experts and the set of active layers are input-dependent.
Early explorations of this combination have appeared in the DeepSeek model family, and the theoretical implications---particularly regarding the interaction between depth-wise and width-wise routing decisions---represent a fertile direction for future work.
The broader vision is a fully adaptive computational graph in which every dimension of the computation (depth, width, precision, even attention span) is conditioned on the input, enabling inference-time compute scaling that allocates resources proportionally to task difficulty.


% ------------------------------------------------------------------------------
\subsection{FFN Sparsity and Post-Training Compression}
\label{subsec:ffn_sparsity}

The observation that large fractions of FFN neurons are inactive for any given input has motivated a rich line of work on exploiting sparsity for inference acceleration, operating at granularities ranging from individual neurons to entire experts.

\paragraph{Contextual sparsity.}
\textbf{Deja Vu}~\cite{liu2023deja} demonstrated that for a given input token, 85--95\% of FFN parameters can be skipped without measurably affecting the output.
Crucially, the set of active neurons varies across inputs---this is contextual sparsity rather than static pruning.
Deja Vu trains a lightweight MLP predictor that, given the previous layer's activation, forecasts which neurons in the current layer will be active, and computes only those neurons.
This approach yields approximately $2\times$ wall-clock speedup on commodity hardware.
The accuracy of the sparsity predictor is the primary bottleneck: overly aggressive prediction causes quality degradation, while conservative prediction limits the achievable speedup.

\paragraph{MoEfication.}
\textbf{MoEfication}~\cite{zhang2022moefication} bridges the gap between dense and sparse architectures by converting a pre-trained dense FFN into a post-hoc MoE structure.
The procedure clusters co-activated neurons into expert groups based on their activation correlation patterns, then trains a lightweight router to select the relevant cluster at inference time.
This enables practitioners to benefit from MoE-style sparsity without the cost and complexity of training an MoE model from scratch, effectively making the dense--MoE boundary a deployment-time rather than design-time decision.

\paragraph{Expert-level pruning.}
In pre-trained MoE models, not all experts contribute equally.
Expert pruning removes underutilised or redundant experts through several criteria: routing frequency (rarely selected experts are pruned), output magnitude (experts with small contributions are discarded), or inter-expert similarity (functionally overlapping experts are merged).
Systematic pruning can reduce the expert count by 30--50\% with minimal quality loss, significantly alleviating the memory footprint problem that plagues MoE inference.

\paragraph{Activation-function-driven sparsity.}
As discussed in Section~\ref{subsec:ffn_activation}, \textbf{ReLU$^2$}~\cite{zhang2024relu2} produces hard zeros for over 90\% of activations, in contrast to the soft near-zero values generated by GELU or SiLU.
Hard sparsity is directly exploitable by hardware: sparse matrix multiplication kernels can skip zero entries entirely, whereas soft sparsity requires thresholding heuristics that introduce approximation error.
This observation has motivated the ``train dense, infer sparse'' paradigm---train with SwiGLU for optimal convergence, then fine-tune or distil into a ReLU$^2$ model for efficient deployment---which is gaining traction as a practical inference optimisation strategy.

\paragraph{Interactions with other compression techniques.}
Sparsity and quantisation are complementary: sparse neurons reduce the number of operations, while quantisation reduces the cost of each operation.
Their joint optimisation---for instance, combining 4-bit weight quantisation with 90\% activation sparsity---has the potential for order-of-magnitude inference speedups, but the interaction effects (e.g., whether quantisation noise disproportionately affects the small set of active neurons) remain insufficiently characterised.


% ------------------------------------------------------------------------------
\subsection{FFN Alternatives: Beyond the MLP Paradigm}
\label{subsec:ffn_alternatives}

The multilayer perceptron---fixed activation functions composed with learnable linear projections---has served as the default channel mixer since the inception of neural networks.
While GLU gating and MoE routing have introduced substantial innovations within this paradigm, a smaller but intellectually significant body of work questions whether the MLP is the only viable design for per-token feature transformation.

\paragraph{Kolmogorov-Arnold Networks.}
\textbf{KAN}~\cite{liu2024kan} draws on the Kolmogorov-Arnold representation theorem, which states that any multivariate continuous function on $[0,1]^n$ can be decomposed as
\begin{equation}
\label{eq:kan}
f(x_1, \ldots, x_n) = \sum_{q=0}^{2n} \Phi_q\!\Biggl(\sum_{p=1}^{n} \phi_{q,p}(x_p)\Biggr),
\end{equation}
where $\Phi_q$ and $\phi_{q,p}$ are univariate functions.
KAN operationalises this theorem by parameterising the univariate functions as learnable B-splines placed on the network's edges rather than its nodes, inverting the standard MLP design in which weights are learnable but activations are fixed.
On low-dimensional scientific computing tasks, KAN exhibits superior approximation efficiency and interpretability compared to MLPs of equivalent size.
However, scaling KAN to LLM-scale models faces formidable obstacles: B-spline evaluation is substantially more expensive than matrix multiplication, efficient GPU kernels for KAN are lacking, and whether KAN's advantages persist in high-dimensional, large-batch regimes is unproven.
Emerging variants---FourierKAN, WaveletKAN, and piecewise-linear approximations---seek to reduce the computational gap while preserving the theoretical appeal.

\paragraph{Adaptive-complexity channel mixing.}
\textbf{Duo-LLM}~\cite{duollm2024} does not replace the MLP but augments it with a parallel lightweight pathway, routing each token to the appropriate complexity level.
This design can be viewed as the channel-mixer analogue of Mixture-of-Depths: rather than skipping entire layers, it adjusts the capacity of the per-token transformation.
The broader implication is that future channel mixers may not be monolithic modules with a single computational profile, but instead offer a continuum of processing intensities tailored to input complexity.

\paragraph{Outlook.}
KAN is more likely to find its niche as a complement to the MLP---deployed in specialised sub-modules requiring high-precision function approximation or interpretability---rather than as a wholesale replacement.
Meanwhile, the convergence of MoE, MoD, and adaptive-complexity designs points toward a unified vision of the channel mixer as a dynamically reconfigurable transformation module.
In this vision, the static notion of ``the FFN'' gives way to a per-token, per-layer decision over what function to apply, how much computation to invest, and which parameters to engage---effectively collapsing expert selection, depth selection, and width selection into a single conditional computation framework.
Realising this vision will require co-design of routing algorithms, hardware accelerators, and training curricula that teach models to allocate computation efficiently from the earliest stages of pre-training.
